\scanpage{82}
Donc on \unsure{trouve} (sans condition de
\unsure{lissité} sur $f$) des \uline{hom de Gysin}
\[
\mathbb{R}_!h\bigl(\mathbb{L}h^*F^\bullet\otimes
 T_{X/\deleted{\unsure{Z}}\inserted{Y}}[2d_{X/Y}]\bigr)
 \longrightarrow \mathbb{R}_!g\bigl(\mathbb{L}g^*(F^\bullet)\bigr)
\]
en particulier
\[
\mathbb{R}_!^{i\inserted{+2d_{X/Y}}}
 \bigl(\mathbb{L}h^*(F^\bullet)\otimes T_{X/Y}\bigr)
 \deleted{[2d_{X/Y}]}
 \longrightarrow \mathbb{R}_!^ig\bigl(\mathbb{L}g^*(F^\bullet)\bigr)
\]
et si $f$ est propre, \unsure{globalement}
\[
H^{i\inserted{+2d_{X/Y}}}
 \bigl(X,\mathbb{L}h^*(F^\bullet)\otimes T_{X/Y}\bigr)
 \deleted{[2d_{X/Y}]}
 \mathbin{\unsure{-|\longrightarrow}}
 H^i\bigl(Y,\mathbb{L}g^*(F^\bullet)\bigr)
\]

Nous admettrons que ces hom satisfont
(\unsure{d'abord} dans le contexte 1$^\circ$, et ``par suite'' dans
le contexte 3$^\circ$) à la condition de
transitivité \unsure{qu'on} devine.
\mspage{81}
